\chapter{Coroutines and async}
\label{ch:async}

\begin{aside}
A coroutine pauses where a function would return. \cpp{}20
gave us the syntax; \maya{} integrates it with the event loop
so a TUI can do async work without blocking the screen.
\end{aside}

\section{Tasks}

\maya{}'s task type wraps a \cpp{} coroutine:

\begin{lstlisting}
Task<int> compute() {
    auto a = co_await fetch(...);
    auto b = co_await fetch(...);
    co_return a + b;
}
\end{lstlisting}

The task runs until the first \code{co\_await}, then yields.
The runtime resumes it when the awaited future is ready.

\section{The scheduler}

\maya{}'s coroutines run on a single thread (the main loop)
by default. \code{co\_await} suspends; the loop continues
processing events; the awaited result fires the resume.

For CPU-bound work, \code{co\_await on(thread\_pool)}
switches to a worker thread. Switch back with
\code{co\_await on(main)} before touching signals.

\section{Awaiters}

Anything with the right shape (\code{await\_ready},
\code{await\_suspend}, \code{await\_resume}) can be
awaited. \maya{} provides:

\begin{itemize}[leftmargin=*]
  \item \code{sleep(duration)}: wait for time to pass.
  \item \code{when\_event(filter)}: wait for an event.
  \item \code{io::read(fd, buf)}: wait for I/O.
  \item \code{when\_signal\_changes(sig)}: wait for a
    signal update.
\end{itemize}

\section{Cancellation}

A task can be cancelled. The cancel signal propagates: the
current await throws \code{Cancelled}; cleanup runs via
destructors.

\maya{} provides a cancellation token:

\begin{lstlisting}
auto token = make_cancel_token();
spawn(task(token));
// later:
token.cancel();
\end{lstlisting}

\section{Effect to task bridge}

An async effect spawns a task that gets cancelled on
dependency change:

\begin{lstlisting}
ctx.async_effect([&]() -> Task<void> {
    auto id = ctx.get(user_id);
    auto data = co_await fetch(id);
    ctx.set(user_data, data);
});
\end{lstlisting}

If \code{user\_id} changes mid-fetch, the in-flight task is
cancelled and a new one starts.

\section{Promises}

For non-coroutine code, \code{Promise<T>} is the bridge:

\begin{lstlisting}
auto p = Promise<int>();
std::thread([&]{
    auto x = expensive_work();
    p.set(x);
}).detach();

Task<int> wait() {
    co_return co_await p.future();
}
\end{lstlisting}

The promise can be set from any thread; the future
awaiter resumes on the main thread.

\section{Avoiding the main thread}

Long CPU work must not run on the main thread, or the
render loop stalls. The worker pool (next chapter) is the
escape hatch.

\section{Recap}

\begin{recap}
\begin{itemize}[leftmargin=*]
  \item Tasks wrap \cpp{} coroutines.
  \item Single-thread by default; \code{on(pool)} to
    offload.
  \item Awaiters for sleep, events, I/O, signals.
  \item Cancellation via tokens.
  \item Promises bridge non-coroutine async code.
\end{itemize}
\end{recap}

\section*{Workshop}

\begin{enumerate}[leftmargin=*]
  \item Build an async fetch widget that loads data and
    displays it.
  \item Cancel the fetch on a typed character (search box
    style).
  \item Run a CPU-heavy task on a worker thread; post the
    result back.
\end{enumerate}
